% ביקורת צוות אדום – סבב 1. קומפילציה: lualatex red-team-round1.tex (פעמיים)
\documentclass[11pt,a4paper]{article}
\usepackage[margin=2.2cm]{geometry}
\usepackage[bidi=basic]{babel}
\babelprovide[main,import]{hebrew}
\babelprovide[import]{english}
\babelfont{rm}{David}
\babelfont{sf}{Arial}
\usepackage{setspace}
\setstretch{1.25}
\setlength{\parindent}{0pt}
\setlength{\parskip}{0.5em}
\newcommand{\en}[1]{\foreignlanguage{english}{#1}}
\usepackage{array}
\usepackage{longtable}
\usepackage{xcolor}

\usepackage{fancyhdr}
\pagestyle{fancy}\fancyhf{}\renewcommand{\headrulewidth}{0pt}
\fancyfoot[C]{\hbox{\textdir TLT\thepage}}
\usepackage[hidelinks]{hyperref}
\newcommand{\cb}{\framebox[1.7ex]{\rule{0pt}{1.3ex}}}
\newenvironment{checks}{\begin{itemize}\renewcommand{\labelitemi}{\cb}}{\end{itemize}}
\newcommand{\crit}{\textcolor{red!70!black}{\textbf{קריטי}}}
\newcommand{\imp}{\textcolor{orange!80!black}{\textbf{חשוב}}}
\newcommand{\nice}{\textcolor{olive}{\textbf{שיפור}}}

\begin{document}

\begin{center}
{\LARGE \textbf{ביקורת צוות אדום: סבב 1}}\\[4pt]
{\large הדוח המכין, מצפן ספקים}\\[2pt]
14.11.2026
\end{center}

\section*{איך להשתמש במסמך}
המסמך מחולק לחמישה חלקים, לפי מה שאת צריכה לעשות בכל אחד מהם:
\begin{enumerate}
\item \textbf{מה תוקן} (חלק 1): לעבור ולוודא שאת מסכימה עם התיקון.
\item \textbf{החלטות שלך} (חלק 2): לבחור אפשרות. בלי ההחלטות האלה אי אפשר לסגור את פרקים 2 ו-5.
\item \textbf{מה להשלים ולאמת} (חלק 3): פרטים שרק אתן יכולות למלא.
\item \textbf{דרך העבודה} (חלק 4): סדר פעולות קבוע לכל סבב עבודה על הדוח.
\item \textbf{צ'קליסט לפני הגשה} (חלק 5).
\end{enumerate}
בכל מקום יש תיבה (\cb) לסימון. מומלץ להדפיס או לסמן בקובץ.

\textbf{מה נבדק:} כל הדוח, מהתקציר ועד פרק 8, משלוש נקודות מבט: המנחה האקדמי (מבנה ועמידה בדרישות), מתודולוג (תוקף התיקוף) ונציג תע"א (שימושיות ואמינות). בבדיקה אוטומטית נמצא שכל 37 המקורות מצוטטים ושאין הפניות שבורות.

% ======================================================
\section{מה תוקן בסבב הזה}
אלה טעויות חד-משמעיות, ולכן תיקנתי אותן ישירות בקבצים. התפקיד שלך: לקרוא את הקטע המתוקן ולוודא שהוא נכון ומובן לך, כי תצטרכי להסביר אותו בהגנה.

{\small
\begin{longtable}{|p{0.5cm}|>{\raggedright\arraybackslash}p{2.2cm}|>{\raggedright\arraybackslash}p{5.6cm}|>{\raggedright\arraybackslash}p{5.3cm}|>{\centering\arraybackslash}p{1cm}|}
\hline
\textbf{\#} & \textbf{איפה} & \textbf{מה הייתה הבעיה} & \textbf{מה תוקן} & \textbf{עברתי} \\ \hline
\endhead
1 & פרק 5, כללי ההחלטה & \textbf{באג לוגי.} הכלל "מידע לא מספיק" רץ לפני הכלל של אירוע בחומרה 2. חברה פרטית כמו \en{Incora}, שעליה מעט מידע אבל ידוע שפשטה רגל, הייתה מקבלת "מידע לא מספיק" ולא אזהרה. & סדר חדש: כללי עצירה, כלל החלטה, ובסוף כלל תקרה. כשל ידוע מוצג בראש הדוח גם כשחסר מידע. נוסף טיפול במצב שאין אף ממד עם מספיק מידע (אחרת חלוקה באפס). & \cb \\ \hline
2 & פרק 5, ניתוח רגישות & הטענה "חלופה ב' עוקפת את ג' רק אם 'מידע חסר' מאבד את רוב משקלו" שגויה מתמטית: גם במשקל 0, ג' מובילה. & חושבה נקודת המפנה האמיתית: ב' עוקפת את ג' רק כשמשקל הישימות עולה על כ-0.45. & \cb \\ \hline
3 & פרק 4, גרף פיזור & ל-\en{Interos} סומן רוחב 9, אבל סכום הציונים שלו ברדאר הוא 7. & תוקן ל-7. & \cb \\ \hline
4 & פרק 2, \en{REQ-FUN-02} & הדרישה "ציון ל-6 מתוך 6 ממדים" סתרה את החלוקה מחדש של משקל כשחסר מידע. & לכל ממד: ציון \textbf{או} סימון "אין מספיק מידע". & \cb \\ \hline
5 & פרק 1, מבנה הדוח & מספרי פרקים נכתבו ידנית. & הוחלפו בהפניות אוטומטיות. & \cb \\ \hline
6 & תקציר & 276 מילים, והמרצה ביקש כ-200. & קוצר לכ-200 מילים. & \cb \\ \hline
\end{longtable}
}

% ======================================================
\section{החלטות שאת צריכה לקבל}
לכל החלטה יש: מה הבעיה, אילו אפשרויות יש, מה ההמלצה שלי, ושאלה שתעזור לך להחליט. \textbf{ההמלצה היא הצעה בלבד}. אתן מכירות את המגבלות שלכן טוב יותר.

\subsection*{החלטה 1 \crit: כיול ובדיקה על אותם מקרים}
\textbf{הבעיה.} בפרק 5 כתוב שהספים (70, 50, 0.4) ייקבעו בבדיקה לאחור. אבל אותם מקרים משמשים גם כדי להוכיח את \en{H1}. אם מכיילים ספים על 10 מקרים ואז בודקים כמה מהם זוהו, התוצאה מנופחת. מנחה מנוסה יזהה את זה.

\textbf{אפשרויות:}
\begin{itemize}
\item \textbf{א. להקפיא את הספים מראש.} קובעים אותם לפני הבדיקה, לפי היגיון וספרות, ואחר כך מציגים ניתוח רגישות: מה היה קורה בספים אחרים. מתאים למדגם קטן.
\item \textbf{ב. לפצל את המקרים.} חלק לכיול, וחלק לבדיקה בלבד. נכון יותר סטטיסטית, אבל דורש מספיק מקרים.
\end{itemize}
\textbf{ההמלצה שלי:} א', אלא אם יימצאו יותר מ-20 מקרים.

\textbf{שאלה מנחה:} כמה מקרי כשל עם ציר זמן מתועד את מעריכה שתמצאו בפועל? אם פחות מ-10, האם פיצול בכלל אפשרי?

\textbf{משפיע על:} פרק 5 (סעיף כללי ההחלטה), פרק 6 (\en{T4.1}), פרק 7 (\en{VAL-02}).

\begin{checks}
\item ההחלטה שלי: \underline{\hspace{8cm}}
\end{checks}

\subsection*{החלטה 2 \crit: קו הבסיס של \en{H1} בחברות פרטיות}
\textbf{הבעיה.} \en{H1} משווה את הציון המשולב מול "בדיקה פיננסית בלבד". אבל לחברה פרטית אין דוחות כספיים, אז קו הבסיס לא מחזיר עליה כלום.

\textbf{אפשרויות:}
\begin{itemize}
\item \textbf{א.} כשאין דוחות, קו הבסיס מחזיר "לא ידוע", וזה נספר כהחמצה. כך בדיוק נמדד הערך המוסף של המערכת בחברות פרטיות.
\item \textbf{ב.} לבדוק את \en{H1} רק על חברות ציבוריות, ואת החברות הפרטיות לבדוק רק ב-\en{H2}.
\end{itemize}
\textbf{ההמלצה שלי:} א', כי היא מראה את הפער שהעבודה באה לסגור. צריך לכתוב את זה במפורש בטבלת ההשערות.

\textbf{שאלה מנחה:} כשאנליסט בתע"א לא מוצא דוחות כספיים, מה הוא עושה היום בפועל? זה מה שקו הבסיס צריך לייצג.

\textbf{משפיע על:} פרק 2 (טבלת ההשערות), פרק 5 (מאגר התיקוף).

\begin{checks}
\item ההחלטה שלי: \underline{\hspace{8cm}}
\end{checks}

\subsection*{החלטה 3 \crit: מה בודקים במקרי הסנקציות}
\textbf{הבעיה.} השער בודק אם הישות \textbf{כבר} ברשימה. ישות שכבר נמצאת שם תיתפס תמיד, ולכן המבחן לא מוכיח הרבה.

\textbf{אפשרויות:}
\begin{itemize}
\item \textbf{א.} להשתמש במקרים האלה רק כבדיקה טכנית שהשער עובד (בדיקת תוכנה, לא תיקוף).
\item \textbf{ב.} להריץ לפי נקודת זמן, \textbf{לפני} ההוספה לרשימה, ולבדוק אם היו סימנים מוקדמים: בעלות, מדינה, חדשות.
\end{itemize}
\textbf{ההמלצה שלי:} שתיהן, אבל לדווח עליהן בנפרד. ב' היא הבדיקה המעניינת מחקרית.

\textbf{שאלה מנחה:} לתע"א חשוב יותר לדעת שהשער אמין, או לקבל אזהרה מוקדמת על ספק שעוד לא נכנס לרשימה?

\begin{checks}
\item ההחלטה שלי: \underline{\hspace{8cm}}
\end{checks}

\subsection*{החלטה 4 \imp: מי ממלא את שאלון \en{H3}}
\textbf{הבעיה.} אחד משני המעריכים הוא המנחה האקדמי, שגם נותן את הציון על הפרויקט. בנוסף, ממוצע של שני מעריכים הוא נתון חלש מאוד.

\textbf{אפשרות מומלצת:} לנסות להשיג 3 עד 5 אנשי רכש או פיתוח עסקי, גם מחוץ לתע"א, ולהשאיר את המנחה כגיבוי בלבד.

\textbf{שאלה מנחה:} את מי אתן מכירות, או דרך מי אפשר להגיע, שעובד ברכש או בהערכת ספקים?

\begin{checks}
\item ההחלטה שלי: \underline{\hspace{8cm}}
\end{checks}

\subsection*{החלטה 5 \imp: ציון "מידע חסר" לחלופה א' במטריצה}
\textbf{הבעיה.} חלופה א' (ידני) קיבלה 2 בטיפול במידע חסר. אבל אנליסט שממלא תבנית קבועה רואה בדיוק אילו שדות ריקים, כך שאפשר לטעון שמגיע לה 3 או 4. מבקר עלול לטעון שהמטריצה "נכתבה אחורה".

\textbf{מה בדקתי:} גם בציון 4, חלופה ג' מנצחת (4.40 מול 3.40).

\textbf{אפשרות מומלצת:} להעלות ל-3 או 4 עם נימוק, ולהציג את הבדיקה הזו כחלק מניתוח הרגישות. זה מחזק את הטיעון.

\begin{checks}
\item ההחלטה שלי: \underline{\hspace{8cm}}
\end{checks}

\subsection*{נושאים חשובים נוספים (לא דורשים בחירה, אבל דורשים עבודה)}
\begin{checks}
\item \imp\ \textbf{מקורות מסחריים וישנים בפרק 2.} הנתונים על משך הבדיקה (\en{HICX}) ועל הסיכון הפיננסי (\en{RapidRatings}) מגיעים מחברות שמוכרות פתרון לבעיה. \en{HICX} ו-\en{ENISA} הם מ-2021. לחפש נתון עדכני ממקור לא מסחרי, או לכתוב במפורש שזה מקור מסחרי.
\item \imp\ \textbf{להתחיל את איסוף מקרי התיקוף כבר עכשיו}, ולא בשבוע 1 של סמסטר ב'. זה הסיכון הגבוה ביותר בפרויקט (\en{DAT-01}).
\item \imp\ \textbf{להוסיף סיכון לוח זמנים לפרק 7.} הנתיב הקריטי הוא בדיוק 14 שבועות, ובניית שישה סוכנים בשלושה שבועות היא אגרסיבית.
\item \nice\ \textbf{פרק 4:} להוסיף משפט שמסביר למה \en{TAM}, \en{SAM} ו-\en{SOM} מוצגים ביחידות שונות (דולרים, מספר ארגונים, נוסחה).
\item \nice\ \textbf{מקורות לפני 2023:} \en{ho2010}, \en{ho2015}, \en{khandpur2020}, \en{caldara2022}, \en{nist2022}. להחליט אילו נחשבים קלאסיים, ולמצוא חלופות עדכניות לשאר.
\item \nice\ \textbf{אורך:} 47 עמודים מול יעד של כ-40.
\end{checks}

% ======================================================
\section{מה להשלים ולאמת}
פרטים שרק אתן יכולות למלא, או מקורות שלא הצלחתי לאמת עד הסוף.

\subsection*{פרטים להשלמה}
\begin{checks}
\item שמות, מיילים ותאריך בעמוד השער.
\item שם ותפקיד של המנחה התעשייתי.
\item חלוקת \en{RACI} בפרק 6: מי סטודנטית א' ומי סטודנטית ב', והאם החלוקה מתאימה לכן.
\item ציוני מטריצת ההחלטה בפרק 5: לעבור על כל ציון ולוודא שאתן יכולות לנמק אותו.
\item ערכי $P$ ו-$I$ בפרק 7: אותו דבר.
\item מילוי \en{AHP} (משקלות סופיים), עם נימוק מהספרות לכל השוואה.
\item 5 עד 10 ספקים יציבים לקבוצת הביקורת, ועוד מקרי כשל.
\end{checks}

\subsection*{מקורות לאימות}
\begin{checks}
\item \en{Prewave}: ההודעה על דירוג \en{Gartner} (הדף לא נטען בבדיקה).
\item \en{Federal Register 2025-19508} (\en{BIS}): הדף חסום לבוטים, צריך לפתוח בדפדפן.
\item \en{Zhao, Ouenniche \& De Smedt}: מספר הכרך.
\item \en{Gani} ועמיתיו: ראשי התיבות של המחברים.
\item מאגר \en{OASIS}: מה חשבון ציבורי מאפשר לראות בפועל.
\item למצוא סקירת \en{MCDM} שפיטה מ-2023 ואילך על בחירת ספקים.
\end{checks}

% ======================================================
\section{דרך עבודה מסודרת}
סדר פעולות קבוע לכל סבב עבודה. המטרה: לא לאבד שינויים, ולא לגלות טעות רק בהגשה.

\subsection*{א. סדר העבודה הכללי (מומלץ)}
\begin{enumerate}
\item \textbf{קודם החלטות, אחר כך כתיבה.} לקבל את החלטות 1 עד 3 (חלק 2). הן משנות את פרקים 2 ו-5, ואין טעם ללטש אותם לפני כן.
\item \textbf{לעדכן את פרקים 2 ו-5} לפי ההחלטות.
\item \textbf{להשלים פרטים} (חלק 3): עמוד שער, \en{RACI}, \en{AHP}.
\item \textbf{לחזק מקורות}: להחליף מקורות ישנים או מסחריים, ולאמת את הרשימה.
\item \textbf{לקצר ל-40 עמודים}: רק בסוף, כשהתוכן יציב.
\item \textbf{סבב ביקורת 2}, ממוקד בפרקים ששונו.
\item \textbf{צ'קליסט לפני הגשה} (חלק 5).
\end{enumerate}

\subsection*{ב. מחזור עבודה על פרק אחד}
בכל פעם שעובדים על פרק, חוזרים על חמשת הצעדים:
\begin{enumerate}
\item \textbf{קוראים} את הפרק כולו לפני שמשנים משהו, ושואלים: מה השאלה שהפרק עונה עליה?
\item \textbf{משנים} בקובץ \en{.tex} של הפרק, בתיקייה שלו.
\item \textbf{מקמפלים}: מתוך התיקייה "דוח מכין - קובץ ראשי" מריצים \en{lualatex main.tex} פעמיים.
\item \textbf{בודקים ב-\en{PDF}}: שהטקסט באנגלית לא הפוך, שטבלאות לא חורגות מהשוליים, שאין "??" במקום מספר פרק או מקור.
\item \textbf{מתעדים}: מה שונה ולמה, במסמך הממצאים של הסבב או ב-\en{CLAUDE.md}.
\end{enumerate}

\subsection*{ג. כשמוסיפים טענה או מספר}
\begin{checks}
\item יש לטענה מקור שאפשר לאמת (\en{DOI}, קישור, דוח רשמי)?
\item המקור נבדק בעצמו, ולא רק דרך מאמר אחר שמצטט אותו?
\item המקור נוסף ל"מקורות - רשימה מאומתת.md" ול-\en{references.tex}, לפי סדר ההופעה?
\item ברור בטקסט אם זו \textbf{עובדה}, \textbf{הנחה} או \textbf{הערכה של הצוות}?
\end{checks}

\subsection*{ד. מלכודות \en{LaTeX} שכבר נתקלנו בהן}
\begin{itemize}
\item טקסט באנגלית בתוך עברית: תמיד עוטפים בפקודה \en{\texttt{\textbackslash en\{...\}}}, אחרת האותיות מתהפכות.
\item טווח מספרים (למשל \en{1--2}) עוטפים באותה פקודה, אחרת סדר המספרים מתהפך.
\item טבלאות: רק \en{longtable}, אחרת סדר העמודות מתהפך.
\item שורה בטבלה שמתחילה ב-[ צריך לעטוף בסוגריים מסולסלים.
\item לא להשתמש בחבילה \en{xurl}: היא תוקעת את הקומפילציה.
\end{itemize}

\subsection*{ה. סבב ביקורת צוות אדום}
אחרי כל שינוי גדול, ולפחות פעם אחת לפני ההגשה:
\begin{enumerate}
\item לקרוא את הדוח מנקודת מבט של \textbf{המנחה}: האם כל דרישה של המדריך ושל שיעור 5 מתקיימת?
\item מנקודת מבט של \textbf{מתודולוג}: האם אפשר להפריך את ההשערות? האם התיקוף הוגן?
\item מנקודת מבט של \textbf{נציג תע"א}: האם הייתי סומך על ההמלצה? מה חסר לי כדי להחליט?
\item לרשום את הממצאים במסמך "סבב \en{N}" בתיקייה הזו, ולחלק אותם ל"תוקן" ול"דורש החלטה".
\end{enumerate}

% ======================================================
\section{צ'קליסט לפני הגשה}
\begin{checks}
\item עמוד שער מלא, בלי מספרי תעודת זהות (המאגר ב-\en{GitHub} עשוי להיות ציבורי).
\item תקציר של כ-200 מילים עם מילות מפתח.
\item תוכן עניינים אוטומטי ומעודכן (קומפילציה פעמיים).
\item מספור עמודים רגיל (1, 2, 3) לאורך כל הדוח.
\item אין "??" בשום מקום.
\item כל טענה, מספר, נוסחה וגרף עם מקור או הסבר מאיפה הם באים.
\item אין מקור שמסומן "לא אומת" או "לאמת".
\item מקורות מלפני 2023 רק אם הם עבודות קלאסיות.
\item הדאטה שאיתו עובדים מוצג במפורש (פרק 5).
\item גאנט עם נתיב קריטי (פרק 6).
\item כל ההחלטות בחלק 2 התקבלו והוטמעו.
\item כ-40 עמודים.
\item הפרויקט עלה ל-\en{GitHub}, עם תיקייה לכל חלק.
\end{checks}

\end{document}
